\documentclass[a4paper]{article}

% 抑制警告
\usepackage{xpatch}

% 中文支持与字体选择
\usepackage[no-math]{fontspec}

\ExplSyntaxOn
\xpatchcmd \fontspec_new_script:nn
  { \__fontspec_warning:nxx }
  { \__fontspec_info:nxx }
  {}{\fail}
\newfontscript{CJK}{none}
\ExplSyntaxOff

\usepackage[fontset = fandol]{ctex}
\usepackage{roboto}

% 数学支持
\usepackage[tbtags]{amsmath}
\usepackage{amsfonts}
\usepackage{amssymb}
\usepackage[warnings-off={mathtools-colon,mathtools-overbracket}, math-style=ISO, mathrm=sym]{unicode-math}
\usepackage{mathtools}
\usepackage{physics}
\usepackage[per-mode = symbol]{siunitx}

% 绘制三线表与跨页表格
\usepackage{booktabs}
\usepackage{longtable}
\usepackage{tabularx}

% 页边距
\usepackage{geometry}

% 页眉与页脚
\usepackage{fancyhdr}

% 图片支持
\usepackage{graphicx}

% 颜色支持
\usepackage{color}
\usepackage[dvipsnames, svgnames, x11names]{xcolor}

% 超链接与书签支持
\usepackage[hidelinks]{hyperref}

% 绘图包
\usepackage[all, pdf]{xy}
\usepackage{tikz}
\usepackage{pgfplots}
\usepackage{circuitikz}
\usepackage{pstricks-add}

% 目录与标题
\usepackage{titlesec}
\usepackage{titletoc}
\usepackage{caption}

% 参考文献
\usepackage[comma, square, super]{natbib}

% Misc
\usepackage{import}
\usepackage{verbatim}

\setmainfont[
    BoldFont       = STIXTwoText-Bold.otf,
    ItalicFont     = STIXTwoText-Italic.otf,
    BoldItalicFont = STIXTwoText-BoldItalic.otf
    ]{STIXTwoText-Regular.otf}
\setmathfont{STIXTwoMath-Regular.otf}

% 页边距与页眉设置
\geometry{a4paper, left = 3.17cm, right = 3.17cm, top = 2.54cm, bottom = 2.54cm}
\setlength{\headheight}{32pt}
\setlength{\headsep}{10pt}
\fancyhead[C]{\includegraphics[scale = 0.06]{res/logo_title.eps}}
\fancyfoot[C]{}
\fancyhead[LR]{}
\fancyfoot[LR]{}
\pagestyle{fancy}

% 标题设置
\titleformat{\section}[block]{\centering\zihao{3}\sffamily}{\thesection\quad}{0em}{}[]
\titleformat{\subsection}[block]{\zihao{4}\sffamily}{\thesubsection\quad}{0em}{}[]
\titleformat{\subsubsection}[block]{\zihao{-4}\sffamily}{\thesubsubsection\quad}{0em}{}[]

\titlecontents{section}[28bp]{\zihao{4}}{\contentslabel{28bp}}{\hspace*{-28bp}}{\titlerule*[0.5pc]{.}\contentspage}
\titlecontents{subsection}[40bp]{\zihao{-4}}{\contentslabel{20bp}}{}{\titlerule*[0.5pc]{.}\contentspage}
\titlecontents{subsubsection}[54bp]{\kaishu\zihao{-4}}{\contentslabel{28bp}}{}{\titlerule*[0.5pc]{.}\contentspage}

\captionsetup{font = sf}

% 行距设置
\linespread{1.5}

% 计数器
\newcounter{code}[section]
\numberwithin{figure}{section}
\numberwithin{table}{section}
\numberwithin{equation}{section}

\renewcommand\thecode{\thesection-\arabic{code}}
\renewcommand\thefigure{\thesection-\arabic{figure}}
\renewcommand\thetable{\thesection-\arabic{table}}
\renewcommand\theequation{\thesection-\arabic{equation}}

% 代码环境配置
\newcommand\inputcode[2]{
    \linespread{1}
    \refstepcounter{code}
    \label{#1}
    \begin{figure}[htbp!]
        \rule[0.5ex]{\textwidth}{0.08em}
        \noindent{\sffamily{\bfseries 代码片段\thecode：}#1} \\
        \rule[0.5ex]{\textwidth}{0.05em}
        \verbatiminput{#2}
        \rule[0.5ex]{\textwidth}{0.08em}
    \end{figure}
    \linespread{1.5}
}

% Misc
\pgfplotsset{compat=1.8}
\newcolumntype{Y}{>{\centering\arraybackslash}X}

% 解决字体基线过低的问题
\makeatletter
\let\original@CJKsymbol\CJKsymbol
\let\original@CJKpunctsymbol\CJKpunctsymbol
\edef\CJKmovesymbol#1{\raise.05em\hbox{\original@CJKsymbol{#1}}}
\edef\CJKmovepunctsymbol#1{\raise.05em\hbox{\original@CJKpunctsymbol{#1}}}
% Only shift non-puncts because puncts seems in their place
\def\CJKraisebaseline{
    \let\CJKsymbol\CJKmovesymbol
}
\def\CJKresetbaseline{
    \let\CJKsymbol\original@CJKsymbol
}
\makeatother

% Activate!
\CJKraisebaseline

\bibliographystyle{gbt7714-numerical.bst}